% Source: https://artofproblemsolving.com/community/c3h1425714_newtons_little_forumla
% Stable identity: aops-topic-1425714-q01
\textbf{(a)} Let $E$ be the shift $a_n\mapsto a_{n+1}$ and $\Delta=E-1$ the difference operator, so $\Delta a_n=a_{n+1}-a_n$. Then $a_{n}=E^{n-1}a_1=(1+\Delta)^{n-1}a_1$, and expanding by the binomial theorem,
\[
a_n=\sum_{j=0}^{n-1}\dbinom{n-1}{j}\Delta^ja_1 .
\]
Here $\Delta^ja_1=\Delta_j$ is the first entry of the $j$th difference row. If the $k$th differences are constant, then all differences of order greater than $k$ are $0$, so the sum stops at $j=k$. (The binomial expansion is valid because $E$ and $1$ commute; equivalently, one can prove the formula by induction on $n$ using Pascal's rule $\binom{n}{j}=\binom{n-1}{j}+\binom{n-1}{j-1}$.)

\textbf{(b)} The difference table is
\[
\begin{array}{cccccccccccc}
1&&5&&14&&30&&55&&91\\
&4&&9&&16&&25&&36\\
&&5&&7&&9&&11\\
&&&2&&2&&2
\end{array}
\]
so $a_1=1$, $\Delta_1=4$, $\Delta_2=5$, $\Delta_3=2$ and
\[
a_n=1+4(n-1)+5\cdot\dfrac{(n-1)(n-2)}{2}+2\cdot\dfrac{(n-1)(n-2)(n-3)}{6}.
\]
Expanding,
\[
a_n=\dfrac{2n^3+3n^2+n}{6}=\boxed{\dfrac{n(n+1)(2n+1)}{6}},
\]
which is the sum of the first $n$ squares, as the second differences $9-4, 16-9,\dots$ suggest.
